%============================================================
% CH2 · PATTERN 1 — EXERCISES, HINTS, ANSWERS
%============================================================
\subsection{Pattern 1: Exercises and Going Further}
\label{ssec:ch2-p1-ex}

\begin{patternbox}{Pattern 1 Summary Card}
\renewcommand{\arraystretch}{1.25}
\begin{tabularx}{\linewidth}{@{}>{\bfseries}l X@{}}
Recognise & ``round trip'', ``loop'', ``deadlock'', ``contradiction in the order'', ``is it a tree / exactly one cycle?'' \\
Directed & colours $0/1/2$; an edge to a \textbf{gray} vertex closes a cycle; black targets are harmless \\
Undirected & skip the edge you came in by (the parent vertex is enough in simple graphs); any other visited neighbour is an ancestor \\
Printing & parent walk from $u$ up to $v$; reverse it for directed graphs; close the loop if the format asks for it \\
Counting & undirected yes/no: cycle iff $m > n - c$; one cycle and connected iff $m = n$ \\
Top traps & ``visited'' instead of gray; no parent skip; start only from vertex 1; recursion on $10^5$ vertices \\
\end{tabularx}
\end{patternbox}

\exlegend
\begin{enumerate}[leftmargin=*, label=\textbf{E1.\arabic*}, itemsep=6pt]
    \item \exeasy\ \pconcept \exown{The stack is the cycle.} Rewrite the directed cycle finder without a parent array: keep the current path in a vector \texttt{path} and an index \texttt{pos[v]} (position of $v$ in \texttt{path}, or $-1$). When an edge reaches a vertex with $\texttt{pos}[v] \ne -1$, cut the cycle out of \texttt{path}. Test it on $D$ (Part~A) and on a graph with a 2-cycle.
    \item \exeasy\ \pclassic \cf{1500}\ \exprob{Cthulhu}{Codeforces 103B}{https://codeforces.com/problemset/problem/103/B}. An undirected simple graph ($n \le 100$) is ``Cthulhu'' if it consists of three or more rooted trees whose roots are joined by a simple cycle. Decide whether the graph is Cthulhu.
    \item \exmedium\ \pclassic \cf{1500}\ \exprob{Fox And Two Dots}{Codeforces 510B}{https://codeforces.com/problemset/problem/510/B}. An $n \times m$ board ($n, m \le 50$) of coloured dots. Is there a cycle of at least 4 distinct, side-adjacent dots of the same colour?
    \item \exmedium\ \ptricky \cf{1600}\ \exprob{New Reform}{Codeforces 659E}{https://codeforces.com/problemset/problem/659/E}. Orient every road of an undirected simple graph ($n, m \le 10^5$) so that as few cities as possible have no incoming road. Print that minimum.
    \item \exhard\ \precent \cf{1700}\ \exprob{Chat Screenshots}{Codeforces 1931F}{https://codeforces.com/problemset/problem/1931/F}. $n$ chat members are ordered by activity, but everyone sees himself on top. $k$ members post screenshots of the order they see. Is there one order consistent with all screenshots? ($n \cdot k \le 2\cdot10^5$.)
    \item \exhard\ \ptricky \exprob{Redundant Connection II}{LeetCode 685}{https://leetcode.com/problems/redundant-connection-ii/}\ \lc{Hard}. A rooted tree on $n \le 1000$ vertices received one extra directed edge. Return an edge whose removal leaves a rooted tree; if several work, the one that appears last in the input.
\end{enumerate}

\subsubsection*{Hints}
\begin{enumerate}[leftmargin=*, label=\textbf{E1.\arabic*}]
    \item \texttt{pos[v] != -1} is exactly ``$v$ is gray''. Reset \texttt{pos[u] = -1} when $u$ is left.
    \item ``Trees hanging from a cycle'' is a connected graph with exactly one cycle. Count instead of searching.
    \item Each colour is a separate grid graph. A grid graph is bipartite, so its cycles are automatically even --- how short can one be?
    \item In a component with a cycle, can every city get an incoming road? In a tree, how many roads are there per city?
    \item A screenshot $a_1, a_2, \dots, a_n$ says nothing about its author $a_1$, but fixes the order of $a_2, \dots, a_n$.
    \item Either some vertex has two parents, or there is a directed cycle, or both. Try deleting the later of the two parent edges and test what remains.
\end{enumerate}

\subsubsection*{Answers}

\begin{solutionbox}[E1.1 The stack is the cycle]
The path vector replaces both the gray colour and the parent array: the cycle is simply the tail of \texttt{path} starting at \texttt{pos[v]}, already in edge direction --- no reversal needed.
\begin{lstlisting}[style=cpp]
#include <bits/stdc++.h>
using namespace std;

// E1.1: the stack IS the cycle. Keep the current path in a vector; pos[v] = index of v in it,
// or -1 if v is not on the path. "Gray" is now simply pos[v] != -1.
vector<vector<int>> adj;
vector<int> path, pos;
vector<bool> done;                                // black
vector<int> cycleFound;

bool dfs(int u) {
    pos[u] = path.size();
    path.push_back(u);
    for (int v : adj[u]) {
        if (pos[v] != -1) {                       // v is on the path: cut the cycle out of it
            cycleFound.assign(path.begin() + pos[v], path.end());
            return true;
        }
        if (!done[v] && dfs(v)) return true;
    }
    path.pop_back();
    pos[u] = -1;
    done[u] = true;
    return false;
}

int main() {
    int n, m;
    cin >> n >> m;
    adj.assign(n + 1, vector<int>());
    for (int i = 0; i < m; i++) {
        int a, b;
        cin >> a >> b;
        adj[a].push_back(b);
    }
    pos.assign(n + 1, -1);
    done.assign(n + 1, false);
    for (int s = 1; s <= n; s++)
        if (!done[s] && dfs(s)) break;
    if (cycleFound.empty()) cout << "acyclic\n";
    else {
        for (int x : cycleFound) cout << x << " ";
        cout << "\n";
    }
}
\end{lstlisting}
On $D$ it prints \textbf{1 2 3}; on \texttt{2 2 / 1 2 / 2 1} it prints \textbf{1 2}; on the DAG \texttt{3 3 / 1 2 / 1 3 / 2 3} it prints \textbf{acyclic}; on \texttt{5 6 / 1 2 / 2 3 / 3 4 / 4 5 / 5 3 / 1 4} it prints \textbf{3 4 5} --- the vertices 1 and 2 before the cycle stay on \texttt{path} but are not part of the answer.
\end{solutionbox}

\begin{solutionbox}[E1.2 Cthulhu]
A connected graph with $n$ vertices has at least $n - 1$ edges; $n - 1$ edges make a tree (no cycle) and every further edge adds one independent cycle. So ``connected with exactly one cycle'' $\iff$ connected and $m = n$. The ``three or more trees'' condition is automatic: in a simple graph a cycle has at least 3 vertices, each the root of its own tree. One search for connectivity, \bigO{n+m}.
\begin{lstlisting}[style=cpp]
#include <bits/stdc++.h>
using namespace std;

// Cthulhu: a connected graph with exactly one cycle  <=>  connected and m == n.
int main() {
    int n, m;
    cin >> n >> m;
    vector<vector<int>> adj(n + 1);
    for (int i = 0; i < m; i++) {
        int x, y;
        cin >> x >> y;
        adj[x].push_back(y);
        adj[y].push_back(x);
    }
    vector<bool> seen(n + 1, false);
    vector<int> st(1, 1);
    seen[1] = true;
    int reached = 1;
    while (!st.empty()) {                       // any search will do: we only count
        int u = st.back(); st.pop_back();
        for (int v : adj[u])
            if (!seen[v]) { seen[v] = true; reached++; st.push_back(v); }
    }
    cout << (reached == n && m == n ? "FHTAGN!" : "NO") << "\n";
}
\end{lstlisting}
Samples: \textbf{FHTAGN!}, \textbf{NO}. Checked against counting all simple cycles explicitly on 400 random graphs with $n \le 6$. The trap: $m = n$ alone is not enough --- a triangle plus a separate edge plus an isolated vertex has $n = m = 5$ but is disconnected.
\end{solutionbox}

\begin{solutionbox}[E1.3 Fox And Two Dots]
Cells of one colour with side adjacencies form an undirected simple graph; run the undirected cycle test of the pattern on it, skipping the cell we came from. A grid graph is bipartite (colour cells like a chessboard), so every cycle has even length; the smallest possible is 4 --- exactly the required minimum, so no length check is needed.
\begin{lstlisting}[style=cpp]
#include <bits/stdc++.h>
using namespace std;

// Fox And Two Dots: undirected cycle detection on the grid graph of each colour.
// A same-coloured visited neighbour that is not the cell we came from closes a cycle.
// Grid graphs are bipartite, so every cycle has even length >= 4 automatically.
int n, m;
vector<string> g;
vector<vector<bool>> seen;
const int dr[4] = {1, -1, 0, 0}, dc[4] = {0, 0, 1, -1};

bool dfs(int r, int c, int pr, int pc) {        // (pr, pc) = the cell we came from
    seen[r][c] = true;
    for (int d = 0; d < 4; d++) {
        int nr = r + dr[d], nc = c + dc[d];
        if (nr < 0 || nr >= n || nc < 0 || nc >= m || g[nr][nc] != g[r][c]) continue;
        if (nr == pr && nc == pc) continue;      // the tree edge back to the parent
        if (seen[nr][nc]) return true;           // back edge: a cycle
        if (dfs(nr, nc, r, c)) return true;
    }
    return false;
}

int main() {
    cin >> n >> m;
    g.resize(n);
    for (auto& row : g) cin >> row;
    seen.assign(n, vector<bool>(m, false));
    for (int r = 0; r < n; r++)
        for (int c = 0; c < m; c++)
            if (!seen[r][c] && dfs(r, c, -1, -1)) { cout << "Yes\n"; return 0; }
    cout << "No\n";
}
\end{lstlisting}
All five samples reproduced (\textbf{Yes, No, Yes, Yes, No}); checked against explicit cycle enumeration on 400 random $2$--$4 \times 2$--$4$ boards with two colours. Recursion is safe here: at most $50 \cdot 50 = 2500$ levels.
\end{solutionbox}

\begin{solutionbox}[E1.4 New Reform]
The model twist: this looks like an orientation problem but is a cycle test per component. A tree component with $k$ cities has $k - 1$ roads, so at least one city gets no incoming road --- and one suffices (orient every road away from any root). A component with a cycle can reach zero: orient the cycle around, then all other roads away from the cycle. So the answer is the \textbf{number of components without a cycle}, i.e.\ with $\#\text{edges} = \#\text{vertices} - 1$.
\begin{lstlisting}[style=cpp]
#include <bits/stdc++.h>
using namespace std;

// New Reform: a component with a cycle can be oriented with no separate city;
// a tree component always keeps exactly one. Answer = number of tree components.
int main() {
    ios::sync_with_stdio(false);
    cin.tie(nullptr);
    int n, m;
    cin >> n >> m;
    vector<vector<int>> adj(n + 1);
    for (int i = 0; i < m; i++) {
        int x, y;
        cin >> x >> y;
        adj[x].push_back(y);
        adj[y].push_back(x);
    }
    vector<bool> seen(n + 1, false);
    int answer = 0;
    for (int s = 1; s <= n; s++) {
        if (seen[s]) continue;
        long long verts = 0, degSum = 0;            // degSum = 2 * (edges of the component)
        vector<int> st(1, s);
        seen[s] = true;
        while (!st.empty()) {
            int u = st.back(); st.pop_back();
            verts++;
            degSum += adj[u].size();
            for (int v : adj[u])
                if (!seen[v]) { seen[v] = true; st.push_back(v); }
        }
        if (degSum / 2 == verts - 1) answer++;      // a tree: no cycle in this component
    }
    cout << answer << "\n";
}
\end{lstlisting}
Samples: \textbf{1, 0, 1}. Checked against trying all $2^m$ orientations on 400 random graphs with $m \le 8$. Isolated cities are tree components too (1 vertex, 0 roads) and count --- they are always separate.
\end{solutionbox}

\begin{solutionbox}[E1.5 Chat Screenshots]
Build an arrow $a_j \to a_{j+1}$ for every pair of consecutive non-authors $j \ge 2$ of every screenshot. A valid order exists iff these arrows contain no directed cycle: if they are acyclic, any topological order works (and puts each author where the others agree); a cycle means two screenshots contradict each other. Iterative colour DFS per test, \bigO{nk}.
\begin{lstlisting}[style=cpp]
#include <bits/stdc++.h>
using namespace std;

// Chat Screenshots: screenshot a1, a2, ..., an says a2 before a3 before ... before an
// (a1 is the author, always shown on top). An order exists iff these arrows have no cycle.
int main() {
    ios::sync_with_stdio(false);
    cin.tie(nullptr);
    int t;
    cin >> t;
    while (t--) {
        int n, k;
        cin >> n >> k;
        vector<vector<int>> adj(n + 1);
        for (int i = 0; i < k; i++) {
            vector<int> a(n);
            for (auto& x : a) cin >> x;
            for (int j = 1; j + 1 < n; j++) adj[a[j]].push_back(a[j + 1]);   // skip the author a[0]
        }
        vector<int> color(n + 1, 0);
        vector<size_t> it(n + 1, 0);
        bool cyclic = false;
        for (int s = 1; s <= n && !cyclic; s++) {
            if (color[s]) continue;
            color[s] = 1;
            vector<int> st(1, s);
            while (!st.empty() && !cyclic) {
                int u = st.back();
                if (it[u] == adj[u].size()) { color[u] = 2; st.pop_back(); continue; }
                int v = adj[u][it[u]++];
                if (color[v] == 1) cyclic = true;            // gray: contradiction
                else if (color[v] == 0) { color[v] = 1; st.push_back(v); }
            }
        }
        cout << (cyclic ? "NO" : "YES") << "\n";
    }
}
\end{lstlisting}
The official sample (ten tests) is reproduced: \textbf{YES YES YES YES NO YES YES YES YES NO}. Checked against trying every order of the members on 1500 random tests with $n \le 5$. Why ignoring the author is right: in a screenshot the author's own position carries no information about the real order.
\end{solutionbox}

\begin{solutionbox}[E1.6 Redundant Connection II]
Every vertex of a rooted tree has exactly one parent, so the extra edge $u \to v$ either gives $v$ a second parent, or points into the root and closes a directed cycle --- or both (it gives $v$ a second parent \emph{and} the old parent edge of $v$ lies on a cycle).
\begin{itemize}[leftmargin=*]
    \item \textbf{No vertex with two parents:} the edges of the unique cycle are the candidates; return the one latest in the input.
    \item \textbf{$v$ has parent edges $A$ (earlier) and $B$ (later):} delete $B$. If no directed cycle remains, $B$ is the answer (it is the later candidate); otherwise $A$ must go.
\end{itemize}
\begin{lstlisting}[style=cpp]
#include <bits/stdc++.h>
using namespace std;

// Redundant Connection II: a rooted tree plus one extra directed edge.
// Case A: some vertex has two parents -> one of its two incoming edges must go.
// Case B: no such vertex -> the extra edge closed a directed cycle; remove its last edge in input order.
class Solution {
    // directed cycle among the given edges (skipping edge 'skip'); returns the edge ids on it
    vector<int> cycleEdges(int n, const vector<vector<int>>& edges, int skip) {
        vector<vector<pair<int, int>>> adj(n + 1);   // (target, edge id)
        for (int i = 0; i < (int)edges.size(); i++)
            if (i != skip) adj[edges[i][0]].push_back({edges[i][1], i});
        vector<int> color(n + 1, 0), parEdge(n + 1, -1);
        vector<size_t> it(n + 1, 0);
        for (int s = 1; s <= n; s++) {
            if (color[s]) continue;
            color[s] = 1;
            vector<int> st(1, s);
            while (!st.empty()) {
                int u = st.back();
                if (it[u] == adj[u].size()) { color[u] = 2; st.pop_back(); continue; }
                int v = adj[u][it[u]].first, id = adj[u][it[u]].second;
                it[u]++;
                if (color[v] == 0) { color[v] = 1; parEdge[v] = id; st.push_back(v); }
                else if (color[v] == 1) {            // back edge id closes the cycle v ... u -> v
                    vector<int> ids(1, id);
                    for (int x = u; x != v; x = edges[parEdge[x]][0]) ids.push_back(parEdge[x]);
                    return ids;
                }
            }
        }
        return {};
    }
public:
    vector<int> findRedundantDirectedConnection(vector<vector<int>>& edges) {
        int n = edges.size();
        vector<int> parentEdge(n + 1, -1);
        int first = -1, second = -1;                 // the two edges into a vertex with 2 parents
        for (int i = 0; i < n; i++) {
            int v = edges[i][1];
            if (parentEdge[v] == -1) parentEdge[v] = i;
            else { first = parentEdge[v]; second = i; }
        }
        if (second == -1) {                          // case B: only a cycle
            vector<int> ids = cycleEdges(n, edges, -1);
            return edges[*max_element(ids.begin(), ids.end())];
        }
        // case A: drop the later edge; if no cycle remains it was the right one
        if (cycleEdges(n, edges, second).empty()) return edges[second];
        return edges[first];
    }
};

int main() {
    Solution s;
    vector<vector<int>> a = {{1, 2}, {1, 3}, {2, 3}}, b = {{1, 2}, {2, 3}, {3, 4}, {4, 1}, {1, 5}};
    vector<int> x = s.findRedundantDirectedConnection(a), y = s.findRedundantDirectedConnection(b);
    cout << x[0] << "," << x[1] << " " << y[0] << "," << y[1] << "\n";   // 2,3 4,1
}
\end{lstlisting}
Examples: \textbf{[2,3]} and \textbf{[4,1]}. Checked against ``try deleting each edge from the last to the first and test for a rooted tree'' on 1500 random trees with one extra edge ($n \le 7$). The trap is the mixed case: when both a second parent and a cycle exist, the answer is the parent edge \emph{on the cycle}, which may be the earlier one even though the problem prefers later edges.
\end{solutionbox}
